\documentclass{situationist-preprint}

\title{More's Utopia Revisited}
\msid{POL-21-b}
\msfamily{POL-21}
\msdiscipline{Political Philosophy}
\mstype{Companion manuscript}
\mskeywords{utopia; Thomas More; reasonable pluralism; perfect society; settlement}

\begin{document}
\maketitle

\begin{abstract}
POL-21-a argues that no society is perfect, because every society settles the trade-offs among the components of the good in a way that some reasonable conception of the good ranks below a feasible alternative. We test the argument against the society that gave utopias their name. Thomas More's Utopia secures abundance, equality and leisure by common property, uniform dwellings and dress, close mutual oversight and a civic religion that excludes unbelievers from office. Each of these is a settlement in the sense of POL-21-a, and for each there are reasonable conceptions of the good that rank an alternative higher. The book itself registers this: its narrator ends by distinguishing what he would wish for from what he expects.
\end{abstract}

\section{Utopia as a settlement}

Definition~2.1 of POL-21-a calls a society perfect only if no reasonable conception of the good gives those who hold it reason to judge some feasible alternative better; Definition~2.2 calls the way a society's institutions and culture resolve the trade-offs among the components of the good a \emph{settlement}; and Theorem~3.1 shows, from Assumptions~2.3 and~2.4, that for every settlement some reasonable conception ranks another higher. More's island is a natural test, because it was designed as the best commonwealth \citep{More1516}.

In Book~II, Raphael Hythloday describes a society without private property, in which every citizen works at a craft and at farming, the working day is limited to six hours, and the surplus is held in common stores. Houses are alike and are exchanged by lot every ten years; dress is uniform; travel requires permission; meals are taken in common halls. Gold and silver are made into chamber pots and the chains of bondmen, so that no one covets them. Religion is tolerated in its variety, but those who deny the immortality of the soul and divine providence are excluded from honors and offices. Grave offenses are punished by bondage.

Each feature resolves a trade-off. Common property and uniform dwellings buy equality and security from want at the cost of privacy and of the expression of individuality through possessions. Mutual oversight buys order at the cost of independence. The civic religion buys a shared basis of trust at the cost of full liberty of conscience. These are exactly the trade-offs among plural goods that Assumption~2.3 of POL-21-a describes.

\begin{proposition}\label{p:utopia}
More's Utopia is not a perfect society in the sense of Definition~2.1 of POL-21-a.
\end{proposition}

\begin{proof}
A conception of the good that gives great weight to personal independence and the cultivation of an individual way of life, of the kind \citet{Mill1859} defended, is reasonable, and its holders have reason to rank above Utopia a feasible society in which property is more widely distributed and oversight lighter, at some cost in equality. A conception that gives weight to liberty of conscience ranks above Utopia a society that does not exclude unbelievers from office. By Definition~2.1 of POL-21-a, Utopia is not perfect.
\end{proof}

The proposition does not depend on the features of Utopia that a modern reader finds repugnant, such as bondage. Remove them, and the settlement of equality against independence remains, and with it the reasonable dissent.

\begin{objection}
More's book is satire, or a thought experiment, not a proposal; it is no test of whether a perfect society is possible.
\end{objection}

\begin{reply}
The interpretive question is open, and we need not settle it. Whatever More intended, the island is described as a working society in which institutions settle trade-offs among goods, and that is all the argument requires. The book's own ending is, moreover, consistent with the conclusion: the narrator grants that there are many features of the Utopian commonwealth he would wish rather than expect to see in our own cities, and he dissents from others, including the abolition of money. A reasonable reader of the best commonwealth disagrees with parts of it \citep{Logan1983}.
\end{reply}

\section{Conclusion}

More's Utopia, the paradigm of the perfect society, is a settlement among plural goods, and reasonable conceptions of the good rank alternatives above it. It thus illustrates the conclusion of POL-21: every utopia conflicts with some reasonable conception of the good. The point is not a criticism of More's island in particular; any improvement that answered its reasonable critics would be a different settlement, open to dissent from another quarter.

\begin{thebibliography}{9}\small
\bibitem[Logan(1983)]{Logan1983} Logan, G. M. (1983). \emph{The Meaning of More's \enquote{Utopia}}. Princeton University Press.
\bibitem[Mill(1859)]{Mill1859} Mill, J. S. (1859). \emph{On Liberty}. John W. Parker and Son.
\bibitem[More(1516)]{More1516} More, T. (1516). \emph{Utopia}. Louvain. Edited by G.~M. Logan and R.~M. Adams, Cambridge University Press, revised edition 2002.
\end{thebibliography}

\end{document}
